\documentclass[runningheads]{llncs}

\usepackage[T1]{fontenc}
\usepackage{array}
\usepackage{booktabs}
\usepackage{graphicx}
\usepackage{longtable}
\usepackage{pgfplots}
\usepackage{pgfplotstable}
\usepackage{url}

\pgfplotsset{compat=1.18}

\title{KNIME Workflows in Scientific Studies: A Semi-Automated Reproducibility Audit}
\titlerunning{KNIME Workflows in Scientific Studies}

\author{Vitalii Kaplan\orcidID{0009-0009-8181-2863}}
\authorrunning{V. Kaplan}
\institute{\email{2333275007@alu.istinye.edu.tr}}

\begin{document}
\maketitle

\begin{abstract}
In this paper, we study how visual workflows are reported in scientific
publications and how incomplete workflow preservation can make older studies
harder to reproduce. Using KNIME as a case study, we develop and apply a
semi-automated evidence pipeline that combines OpenAlex bibliometrics, a
100-record top-cited article audit registry, workflow retrieval, and opening
and execution checks in a current KNIME environment. The OpenAlex search shows
that KNIME is visible across a broad scientific literature, while the article
assessment shows that workflow figures and textual descriptions are more
common than downloadable workflow artifacts. Of the 100 audited records, 22
report downloadable or linked KNIME workflows, and recorded retrieval attempts
obtained workflow artifacts or workflow directories for 18. In the
recorded manual checks, five article records had at least one workflow execute
successfully. These results show that incomplete artifact and environment
preservation limits the long-term reproducibility of KNIME-based studies.

\keywords{Workflow reproducibility \and KNIME \and Bibliometrics \and
Research artifacts}
\end{abstract}

\section{Introduction}

KNIME is a widely used open-source platform for constructing scientific
data-analysis workflows, with more than a decade of use across published
studies \cite{Berthold2009KNIME}. Its visual workflows represent analyses as
connected nodes with
stored settings, data ports, and extension dependencies, which can make
computational work easier to inspect, reuse, and repeat. This representation
provides more operational detail than a narrative method description, but it
does not remove the need to preserve the executable workflow artifact and its
context. Reproduction requires access to the workflow file, KNIME version,
extensions, data, and execution environment. Figures and textual summaries do
not provide sufficient substitutes for these materials.

This issue matters for visual workflow platforms in general, because their
reproducibility value depends on preserving executable workflow artifacts and
the software context needed to interpret them. We study KNIME because it is an
open-source visual workflow platform with a long public development history and
substantial visibility in scientific literature. Our OpenAlex query returned
963 article-type records whose titles or abstracts match \texttt{KNIME}. The
matching records span platform papers, extension papers, tool-comparison
papers, and domain studies that use KNIME as an analysis interface. KNIME is
therefore a suitable case for studying whether visual workflows are preserved
as executable artifacts or only reported through descriptions and figures.

Our top-cited article collection covers the 100 most-cited KNIME-matching
records, with not-assessed placeholders retained where local full text is
unavailable. The structured audit reported in
Table~\ref{tab:top-cited-assessment} shows an uneven publication pattern.
Twenty-two records report downloadable or linked KNIME workflow files, and
workflow artifacts or workflow directories were obtained for 18 article
records in the recorded retrieval attempts. Other links were stale, blocked,
incomplete, or did not lead to a confirmed KNIME workflow artifact. In the
recorded manual execution checks, five article records had at least one
workflow execute successfully. By comparison, 49 records present workflow
screenshots or figures, and 68 describe workflow steps or nodes in text. These
forms of presentation do not preserve the complete executable workflow
context.

These reporting practices create a reproducibility risk because later readers
may lack the executable artifact, input data, recorded KNIME version, or
extension context needed to reconstruct and run the analysis. We therefore
combine OpenAlex bibliometrics, top-cited article retrieval and assessment,
workflow-artifact retrieval, and current-version opening and execution
attempts. The goal is to measure how KNIME workflows are reported and
preserved, and to distinguish publication-level claims of workflow availability
from artifacts that were actually obtained and tested.

The rest of the paper is organized as follows. Section~2 reviews related work
on computational reproducibility, open-source workflow systems, environment
preservation, and research artifacts. Section~3 presents the bibliometric
method and semi-automated evidence pipeline. Section~4 describes the resulting
datasets. Section~5 reports the bibliometric, article-audit,
workflow-retrieval, and manual opening and execution results. Section~6
discusses the implications, limitations, and future extensions, and Section~7
concludes the paper.

\section{Literature Review}

Reproducibility is a central requirement for computational research because
published results increasingly depend on software, data transformations, and
execution environments. Provenance research makes this requirement explicit:
visual analytics systems need records of data, analytical processes, and human
interpretation to make later inspection and reuse possible
\cite{Walker2013ProveML}. Studies of executable research artifacts show the
same problem in practice. Samuel and Mietchen found that biomedical Jupyter
notebooks often fail during dependency installation or re-execution even when
the notebooks themselves are available
\cite{SamuelMietchen2024JupyterReproducibility}.

Open-source scientific software helps reproducibility by making algorithms,
interfaces, and execution infrastructure inspectable. Weka, for example, is
presented as an extensible open-source data-mining environment with graphical
interfaces, reusable configurations, and plugin mechanisms
\cite{Hall2009WEKA}. KNIME is similarly used as an open-source visual workflow
platform, and recent GIScience work explicitly proposes KNIME-based workflows
as an approach for advancing replicable and reproducible research
\cite{Fu2025GIScienceKNIME}. However, openness of the platform alone is
insufficient. A workflow-based study also requires the executable workflow
file, input data, platform version, installed extensions, and other software
dependencies. Container research makes the same point at the environment
level: scientific computations are easier to reproduce when complete software
environments can be preserved and moved between systems
\cite{Kurtzer2017Singularity}. More broadly, the testing literature shows that
executable artifacts can remain sensitive to operating systems, dependency
versions, timing, external services, and other environmental factors
\cite{Rasheed2023FlakinessReview}.

For visual workflow studies, preservation must include more than an image of
the workflow graph. Node settings, extension dependencies, input paths, and
version context are part of the executable method. The distinction between a
reported workflow, a downloadable artifact, an artifact that opens, and one
that executes successfully is therefore central to the present audit.

\section{Methods}

All scripts and data used for this study, except the other articles assessed
as study objects, are available in the public GitHub repository
\cite{KaplanReproducibilityRisksArticle2026}. Relative paths reported in this
paper should be interpreted as paths within that repository.

\subsection{Bibliometric Evidence}

We collected bibliometric evidence from the OpenAlex Works API
\cite{Priem2022OpenAlex} on June 29, 2026. The query used OpenAlex
title-and-abstract search for the word \texttt{KNIME} and restricted the
result set to OpenAlex records whose type is \texttt{article}. The same query
can be reproduced in the
OpenAlex web interface by selecting ``Search title and abstract'', entering
\texttt{KNIME}, and filtering the result type to \texttt{article}. The
corresponding web-interface request can be represented as:

\begin{quote}
\path{https://openalex.org/works}\\
\path{page=1}\\
\path{sort=relevance_score:desc}\\
\path{search.title_and_abstract=KNIME}\\
\path{filter=type:article}
\end{quote}

The first API request, which returns the first 200 records, is reproducible as:

\begin{quote}
\path{https://api.openalex.org/works}\\
\texttt{?search.title\_and\_abstract=KNIME}\\
\texttt{\&filter=type:article}\\
\texttt{\&per-page=200}\\
\texttt{\&cursor=*}
\end{quote}

OpenAlex uses cursor pagination for large result sets. After the first request,
the collector reads \texttt{meta.next\_cursor} from each response and submits
the same request with \texttt{cursor=<next-cursor>} until no further results
are returned. The value \texttt{per-page=200} is the maximum page size used
here.

The collection script is
\path{scripts/collect_openalex_knime_works.py}.
It stores raw and processed collection artifacts as follows:
\begin{itemize}
\item raw OpenAlex response pages:
\path{data/original/openalex/pages/}
\item complete record stream:
\path{data/original/openalex/works.jsonl}
\item compact CSV view:
\path{data/original/openalex/works.csv}
\item endpoint, query parameters, collection time, result count, and page
manifest:
\path{data/original/openalex/metadata.json}
\end{itemize}
The current run collected 963 records across 5 OpenAlex response pages.
Derived bibliometric tables are generated from the JSONL file by
\path{scripts/build_openalex_knime_bibliometrics.py}.

\subsection{Semi-Automated Evidence Pipeline}

For the article-level audit, we developed an open-source, semi-automated
evidence pipeline composed of scripts and JSON configuration files. The
process is specified in \path{scripts/audit_script_chain.json}, which serves as
the workflow index for this part of the study. Each step consumes recorded
outputs from earlier steps and writes an explicit intermediate or final
artifact.
The current chain contains fifteen scripted steps and one manual checkpoint.
Thirteen scripted steps are rule-based. Two scripted steps use an LLM.
We use the LLM only where rule-based extraction is insufficient. The first LLM
step classifies workflow relevance from downloaded reference-page content. The
second LLM step fills supplementary article-level audit flags from the
processed text of the articles. The LLM calls use temperature 0 to reduce
variation. Repeated execution may still depend on the model and service version
used.

The pipeline is largely automated, but it deliberately permits human
intervention at two boundary points: obtaining article PDFs at the beginning
of the process and obtaining workflow artifacts near the end, before formal
result preparation.
Both tasks may involve institutional access, protected publisher pages,
supplement portals, captchas, and changing web interfaces. The process can
therefore benefit from human assistance at the beginning and end, while the
intermediate steps execute automatically from recorded evidence.

The implementation was run with Python 3.14.5 and uses a separate
\path{.venv-playwright} virtual environment for browser-dependent steps.
Article PDF conversion uses GROBID \texttt{0.9.0-crf} through the Docker image
\texttt{grobid/grobid:0.9.0-crf} \cite{GROBIDSoftware2026}. Browser-assisted
fetching uses Playwright for Python 1.61.0 with Chromium/Chrome
\cite{PlaywrightPython2026}. The two LLM-supported steps call the OpenAI Chat
Completions API directly, using the configured model
\texttt{gpt-4.1-mini} and temperature 0 \cite{OpenAIAPI2026}. OpenAlex is the
bibliographic source for article discovery and ranking
\cite{Priem2022OpenAlex}. The exact commands, script order, prompt files, and
intermediate outputs are recorded in \path{scripts/audit_script_chain.json}.

For reproducibility of the LLM-supported steps, API configuration is separated
from prompts and data. The scripts read the API key from \path{.env} or the
process environment as \texttt{OPENAI\_API\_KEY}; this local credentials file
is not part of the research data. The model can be overridden with
\texttt{OPENAI\_MODEL}, but the recorded run used the default
\texttt{gpt-4.1-mini}. Both LLM scripts use the Chat Completions endpoint
\path{https://api.openai.com/v1/chat/completions}, temperature \texttt{0.0},
and at most three retries for transient server-side failures. The prompts are
stored as JSON configuration files:
\path{data/processed/audit/article_reference_classification_prompt.json} for
reference-page classification and
\path{data/processed/audit/article_supplementary_flag_prompt.json} for
article-level supplementary flags. The derived JSON outputs record the script
name, input file, prompt configuration, model, temperature, selection scope,
summary counts, and per-article classifications. This makes the exact prompt
and local evidence auditable, while recognizing that future reruns may differ
if the hosted model implementation changes.

The following paragraphs describe each step in the ordered script chain. For
each step, we identify its purpose, implementation script, principal inputs or
outputs, and any point at which manual intervention may be required.

\textbf{Step 1} collects KNIME-matching article records from the OpenAlex works API with
\path{scripts/collect_openalex_knime_works.py}. The query searches for KNIME in
article titles and abstracts and stores the original API response pages, a
complete JSONL stream, a compact CSV view, and query metadata. This step
preserves the bibliographic source records before any article-level audit
decisions are made.

\textbf{Step 2} builds the bibliometric summaries and the citation-ranked seed table
with \path{scripts/build_openalex_knime_bibliometrics.py}. The script reads the
OpenAlex JSONL export and derives processed CSV and JSON files used to select
the top-cited audit candidates. The citation-ranked seed table is then used as
the article list for downstream download, parsing, URL collection, and audit
report construction.

\textbf{Step 3} attempts to download article PDFs by citation-rank range with
\path{scripts/download_openalex_rank_range_articles.py}. This step is
semi-automated because OpenAlex links and open-access URLs do not always expose
the full paper, and some papers require lawful institutional access. The script
records download attempts and stores retrieved PDFs under
\path{data/original/articles/}; failed or protected records can be completed by
manual download without changing the later audit logic.

\textbf{Step 4} converts the local PDFs into semantic article HTML with
\path{scripts/extract_article_grobid_html.py}. The current parser uses a local
GROBID service to produce TEI XML and then renders an HTML reading copy with
article metadata, abstract, body sections, paragraphs, and bibliography
structure. This replaced earlier plain-text extraction because the downstream
steps need stable article identity, paragraphs, and reference context.

\textbf{Step 5} collects article metadata and article-related URLs with
\path{scripts/collect_article_urls.py}. The script joins the OpenAlex seed
table, the canonical PDF registry, the GROBID TEI files, and the generated HTML
files. At this stage it does not classify URLs or infer workflow availability;
it only records metadata and normalized URLs that can be traced to the article
content.

\textbf{Step 6} fetches the referenced URL pages with
\path{scripts/fetch_article_url_pages.py}. For each unique URL, the script
records HTTP status, redirect chain, final URL, response headers, and a stored
copy of the response body where available. This creates local evidence for
later classification and avoids asking the language model to rely on live web
state.

\textbf{Step 7} attaches the URL-fetch metadata back to the article records with
\path{scripts/attach_url_fetch_metadata.py}. This rule-based step augments
each URL entry with fetch status, HTTP code, final URL, redirect count, and the
path to any stored page. It is metadata attachment only; it does not interpret
whether a reference is a workflow, software page, dataset, article landing
page, or documentation.

\textbf{Step 8} performs a second-stage fetch for reference pages that ordinary HTTP
requests cannot access, using \path{scripts/browser_fetch_403_urls.py}. This
step uses a browser when useful, stores rendered pages, and records challenge
or blocked states explicitly. The process
does not solve captchas or bypass access controls; protected pages are treated
as evidence of access friction rather than forced downloads.

\textbf{Step 9} classifies reference pages for workflow obtainability with
\path{scripts/classify_article_references_with_llm.py}. The language model is
given URL metadata and locally cached reference-page excerpts, not the whole
article as its primary evidence. The purpose is to decide whether a reference
can help a reader obtain a KNIME workflow artifact, whether it only points to
software, documentation, data, or a general page, or whether manual checking is
needed.

\textbf{Step 10} classifies supplementary article-level flags with
\path{scripts/classify_article_supplementary_flags_with_llm.py}. This LLM step
uses the article HTML, metadata, and reference classifications to fill bounded
boolean fields comparable to a manual audit checklist: KNIME use, KNIME version
reporting, workflow screenshots or figures, textual workflow or node
descriptions, input-data availability, code or scripts, and extension or
plugin information. These flags provide article-level context around the
reference-page workflow evidence.

\textbf{Step 11} builds a compact article audit report with
\path{scripts/build_article_audit_report.py}. This rule-based join combines
article metadata, the supplementary flag fields, and the workflow-relevant
reference classifications into
\path{data/processed/audit/article_audit_report.json}. The report is the
central article-level audit artifact used by later workflow-download and table
generation steps.

\textbf{Step 12} attempts automated workflow-artifact downloads from the workflow
references in the audit report with
\path{scripts/download_workflow_artifacts.py}. The script uses rule-based
HTTP retrieval and, where appropriate, browser-assisted discovery to look for
KNIME workflow archives, \texttt{workflow.knime} files, or archives containing
workflow files. It writes attempt status, local directories, downloaded files,
and discovered workflow files to the workflow-reference inventory.

\textbf{Step 13} is the manual workflow-download checkpoint. The auditor inspects
\path{data/processed/audit/knime_downloadable_workflow_references.json} and
tries to download missing workflow artifacts where access is allowed. The
inventory JSON is not edited manually at this stage; the manual action is to
place obtained files in the article-specific directory under
\path{data/original/workflows/}. This preserves a clear boundary between human
retrieval and scripted evidence updating.

\textbf{Step 14} updates the article report from retrieved workflow evidence with
\path{scripts/update_article_report_from_workflow_inventory.py}. The script scans
the workflow inventory and the actual files under \path{data/original/workflows/}
and updates the article audit report when a downloadable KNIME workflow has
been obtained. This makes the report reflect the local evidence rather than
only the earlier URL classification.

\textbf{Step 15} generates the article-audit summary table source with
\path{scripts/build_article_audit_tables.py}. The script reads the compact
article audit report and writes the CSV source for the publication-level flag
summary. The resulting data are therefore
derived from the final structured report rather than maintained by hand.

\textbf{Step 16} generates the KNIME-use workflow-reporting table source with
\path{scripts/build_knime_use_workflow_reporting_table.py}. This step combines
the article-level audit report with the workflow-reference inventory, so that
publication-level workflow-reporting signals and actual obtained workflow
evidence are kept separate but can be summarized consistently.

The resulting evidence pipeline is reproducible at the level of inputs, scripts,
configuration files, cached pages, LLM prompts, and derived JSON reports, while
still acknowledging that publication and workflow retrieval are not fully
automatable. Manual actions are restricted to acquiring protected article PDFs
and workflow artifacts where access is lawful; the subsequent interpretation,
report updating, and table construction are scripted from recorded evidence.

All project Python scripts are intended to be run with Python 3.11 or newer.
The current scripts use only the Python standard library. No third-party
Python packages are required for the OpenAlex collection and rule-based audit
steps. Browser-assisted, GROBID, and LLM-supported steps use the separately
recorded tools and services described above.

\section{Data}

\subsection{Bibliometric Evidence}

The OpenAlex bibliometric data are organized into original API responses and
processed analysis tables. The original data are stored under
\path{data/original/openalex/}.
They include raw response pages, a complete JSONL export, a compact CSV export,
and query metadata. These files preserve the article-level records used to
measure the size, time trend, venues, fields, and accessibility of the
KNIME-matching literature.

Derived bibliometric tables are stored under
\path{data/processed/openalex/}.
Their roles in the analysis are:

\begin{itemize}
\item \path{data/processed/openalex/openalex_knime_articles.csv}: compact article-level
metadata used as the main processed bibliography table.
\item \path{data/processed/openalex/openalex_knime_articles_by_year.csv}: annual publication
counts used to describe the time trend of KNIME-matching literature.
\item \path{data/processed/openalex/openalex_knime_top_venues.csv}: source-level aggregation
used to identify journals, proceedings, repositories, and platforms where the
matching records appear.
\item \path{data/processed/openalex/openalex_knime_top_domains.csv}: OpenAlex domain
aggregation used to describe the broad disciplinary distribution.
\item \path{data/processed/openalex/openalex_knime_top_fields.csv}: OpenAlex field
aggregation used to describe the finer disciplinary distribution.
\item \path{data/processed/openalex/openalex_knime_most_cited.csv}: citation-ranked subset
used to select influential records for later manual assessment.
\item \path{data/processed/openalex/openalex_knime_likely_workflow_platform.csv}: heuristic
subset used to prioritize papers likely to use KNIME as a workflow platform.
\end{itemize}

The processed OpenAlex directory also contains additional summary files:
\begin{itemize}
\item \path{data/processed/openalex/openalex_knime_open_access_availability.json}
\item \path{data/processed/openalex/openalex_knime_role_classification_summary.json}
\end{itemize}
These files summarize artifact-access signals and role-classification counts.

\subsection{Top-Cited Article Assessment}

The top-cited article dataset records the structured assessment of the 100
most-cited KNIME-matching OpenAlex records. The citation-ranked source list is
\path{data/processed/openalex/openalex_knime_most_cited.csv}, and the local
article registry is stored under \path{data/original/articles/}. Local PDFs,
where available, are preserved in that directory, while download-attempt logs
are stored under \path{data/processed/audit/logs/}.

Article text and references are represented by GROBID outputs. Semantic HTML
files are stored under \path{data/processed/articles/}, and TEI XML files are
stored under \path{data/processed/articles/grobid_tei/}. These outputs provide
article text, section structure, bibliography context, and machine-readable
links for the later URL and LLM steps.

The URL extraction step writes
\path{data/processed/audit/article_url_collection.json}. For each article, this
file stores rank, canonical identifiers, PDF path, TEI path, HTML path,
OpenAlex metadata, and article-related URLs extracted from the GROBID outputs.
The fetched reference-page cache is stored under
\path{data/processed/audit/pages/}; browser-rendered fetches for selected
blocked pages are stored under \path{data/processed/audit/browser_pages/}.
These caches include HTTP status, redirects, final URLs, headers, stored body
paths, and fetch errors or challenge pages where applicable.

The LLM-derived article-audit data are split into two files.
\path{data/processed/audit/article_reference_llm_classifications.json}
classifies cached article references by workflow obtainability and records the
page evidence used for those classifications.
\path{data/processed/audit/article_supplementary_llm_flags.json} records
article-level flags for KNIME use, version reporting, workflow figures,
workflow text descriptions, data availability, code availability, and
extension information.

The compact report used for the article-level tables is
\path{data/processed/audit/article_audit_report.json}. It merges article
metadata, supplementary flags, and workflow-relevant reference classifications
into one record per article. Workflow-download attempts and local workflow
paths are recorded in
\path{data/processed/audit/knime_downloadable_workflow_references.json}; manual
KNIME-opening screenshots are stored under each workflow directory's
\path{opened/} subdirectory. Table sources are generated under
\path{article_bibliometric/tables/} from these JSON files rather than edited
manually.

\section{Results}

\subsection{Bibliometric Evidence}

The OpenAlex title-and-abstract query returned 963 article-type records. This
establishes that KNIME is visible in a sizeable scientific literature and that
workflow preservation is relevant beyond platform-maintenance research. The
annual series indicates that KNIME-matching literature appears in OpenAlex
from the mid-2000s and grows into a
sustained publication stream after 2015 (Fig.~\ref{fig:openalex-by-year}). The
highest annual count in the current export is 104 records in 2023. Counts for
2026 are incomplete because the query was collected on June 29, 2026.

\begin{figure}
\centering
\begin{tikzpicture}
\begin{axis}[
  width=0.95\textwidth,
  height=0.34\textwidth,
  ybar,
  bar width=3.5pt,
  xlabel={Publication year},
  ylabel={Article records},
  ymin=0,
  xmin=2005,
  xmax=2027,
  enlarge x limits=0.02,
  xtick={2006,2010,2015,2020,2026},
  ymajorgrids=true,
  grid style={gray!25},
  tick label style={font=\scriptsize},
  label style={font=\small},
]
\addplot table[
  x index=0,
  y index=1,
  col sep=comma,
  ignore chars={\"}
] {tables/openalex_knime_articles_by_year.csv};
\end{axis}
\end{tikzpicture}
\caption{OpenAlex article records matching \texttt{KNIME} in title or
abstract by publication year.}
\label{fig:openalex-by-year}
\end{figure}

The domain distribution shows that the KNIME-matching literature is not
confined to one disciplinary area (Table~\ref{tab:openalex-domains}). OpenAlex
assigns just over half of the records to Physical Sciences, but the remaining
records are spread across Social Sciences, Life Sciences, and Health Sciences.
This matters for the study because workflow reproducibility risk is not only a
software-engineering concern. If KNIME workflows are used across computational
chemistry, biomedical analysis, education, business analytics, and related
areas, then missing workflow artifacts and undocumented dependencies can
affect published computational evidence in several research communities.

\begin{table}
\centering
\caption{OpenAlex domains for KNIME title-and-abstract matches.}
\label{tab:openalex-domains}
\begin{tabular}{lrr}
\toprule
Domain & Articles & Share \\
\midrule
Physical Sciences & 505 & 52.44\% \\
Social Sciences & 187 & 19.42\% \\
Life Sciences & 149 & 15.47\% \\
Health Sciences & 122 & 12.67\% \\
\bottomrule
\end{tabular}
\end{table}

The ten largest OpenAlex fields are led by Computer Science (350),
Biochemistry, Genetics and Molecular Biology (115), Medicine (92),
Engineering (71), and Social Sciences (66). These field counts show why a
paper about KNIME reproducibility cannot focus only on the KNIME platform
literature. The central reproducibility question is how scientific studies
that use KNIME as a tool report the executable workflow, data, versions, and
extensions needed for later reproduction.

Open-access metadata indicate 619 open-access records, 625 records with some
full-text availability, and 421 records with a PDF signal. These signals are
important because later workflow audits require access to paper text,
supplements, repositories, and sometimes PDF-only workflow descriptions. The
counts are only availability indicators. They do not by themselves establish
that a workflow artifact, KNIME version, input data, or extension list is
available.

The heuristic role classification separates likely papers about KNIME or KNIME
extensions from papers that likely use KNIME as an analysis tool. The current
rules classify 156 records as about KNIME or a KNIME extension, 314 as using
KNIME as a tool, 28 as about or mentioning KNIME in the title, and 465 as
ambiguous OpenAlex matches. The same heuristic marks 410 records as likely
using KNIME as a workflow platform. These labels are useful for triage: they
identify a substantial candidate set where workflow reporting can matter, but
final case-study claims require manual full-text and artifact assessment.

\subsection{Top-Cited Article Assessment}

Of the 100 citation-ranked KNIME-matching records, 37 article PDFs were
downloaded by the retrieval script and 42 were obtained manually, giving 79
local PDFs in total. The top-cited audit therefore shows that article access
itself is already a limiting condition for reproducibility assessment. This
means that the automated middle of the pipeline can evaluate most,
but not all, top-cited records from local article text and extracted
references. The remaining records stay in the denominator, because excluding
them would overstate the assessability of the literature.

The URL and reference-page stages add a more precise view of workflow
availability than article text alone. The compact report contains
workflow-relevant linked resources for 22 article records, with 68 linked
resources retained after filtering generic article pages, datasets, software
documentation, and malformed or unrelated URLs. This result is important
methodologically: the clearest evidence for whether a reader can obtain a
workflow can appear on the referenced page itself rather than in the sentence
that cites the URL.

The LLM reference-classification step reduced the URL set to resources that
could plausibly support workflow retrieval and assigned bounded categories for
downloaded artifacts, workflow landing pages, direct workflow links, and
resources requiring manual inspection. The broader workflow inventory retains
31 candidate workflow-reference records because it preserves retrieval
attempts even when the compact article-level flag is stricter. This separation
allows the audit to distinguish evidence that an article or reference suggests
a workflow is available from evidence that a workflow artifact was obtained
locally.

The supplementary LLM flag step provides the publication-level reporting
context for those workflow links. Seventy-five records use KNIME as a workflow,
tool, interface, or implementation context. Of the full 100-record registry, 68
records describe workflow steps or nodes in text, and 49 show workflow
screenshots or figures. In contrast, 18 report a KNIME version, 22 report a
downloadable or linked workflow, 23 report extension dependencies, and four
report an extension installation source. Visual or narrative workflow
presentation is therefore more common than the executable and environment
context needed for later reproduction.

Restricting the reporting signals to the 75 KNIME-use cases changes two of the
all-record counts: 48 of the KNIME-use articles show workflow screenshots or
figures, and 22 report extension or plugin dependencies. The corresponding
all-record counts are 49 and 23 because each signal also occurs once in a
comparative paper that is not classified as a KNIME-use case. The remaining
signals in Table~\ref{tab:knime-use-assessment} are fully contained within the
75 KNIME-use cases.

The workflow-download stage converts reported or linked workflow evidence into
local artifact evidence. Eighteen article records yielded workflow files,
archives, or workflow directories. Four were obtained automatically and
fourteen required manual download. This confirms the semi-automated character
of the boundary step: scripted retrieval can find some artifacts, but protected
pages, stale links, repository layouts, and supplement portals still require
human inspection.

Table~\ref{tab:top-cited-assessment} summarizes the statistics-ready fields in
the compact article audit report. Data and code signals are recorded separately
from workflow-artifact evidence. We do not verify the availability of reported
data or code in this table; we record whether the article and linked evidence
report those resources.

\begingroup
\scriptsize
\renewcommand{\arraystretch}{1.18}
\begin{longtable}{|%
  >{\raggedright\arraybackslash}p{0.34\textwidth}|
  r|
  r|
  >{\raggedright\arraybackslash}p{0.37\textwidth}|}
\caption{Flag summary for the top-cited article audit. Shares use all
100 records in the registry as the denominator.}
\label{tab:top-cited-assessment}\\
\toprule
Audit category or flag & Count & Share of 100 & Interpretation \\
\midrule
\endfirsthead
\toprule
Audit category or flag & Count & Share of 100 & Interpretation \\
\midrule
\endhead
\bottomrule
\endfoot
Audit records & 100 & 100.0\% & Most-cited KNIME-matching article records selected from OpenAlex \\
Records with local PDF & 79 & 79.0\% & Source PDF recorded in the compact audit report \\
\midrule
Uses KNIME & 75 & 75.0\% & KNIME used as a workflow, tool, interface, or implementation context \\
Workflow or nodes described in text & 68 & 68.0\% & Named workflow steps, nodes, modules, or components recorded \\
Workflow screenshots or figures & 49 & 49.0\% & Workflow shown visually in article or supplement \\
Reports KNIME version & 18 & 18.0\% & Specific KNIME version value found \\
Reports downloadable workflows & 22 & 22.0\% & Executable workflow files provided or linked \\
Reports extension/plugin dependencies & 23 & 23.0\% & KNIME extension or node-library dependency reported \\
Reports extension installation source & 4 & 4.0\% & Update site, project URL, or installation source reported \\
\midrule
Provides code or scripts & 20 & 20.0\% & Code repository, scripts, or software code reported \\
Reports input-data availability & 66 & 66.0\% & Data availability stated, with or without direct URL \\
Direct input-data resource & 20 & 20.0\% & Direct data URL or resource pointer recorded \\
\end{longtable}
\endgroup

\begingroup
\scriptsize
\renewcommand{\arraystretch}{1.18}
\begin{longtable}{|%
  >{\raggedright\arraybackslash}p{0.62\textwidth}|
  r|
  r|}
\caption{Workflow-reporting and retrieval signals among the 75 assessed KNIME-use cases.}
\label{tab:knime-use-assessment}\\
\toprule
Audit flag & Count & Share of KNIME-use cases \\
\midrule
\endfirsthead
\toprule
Audit flag & Count & Share of KNIME-use cases \\
\midrule
\endhead
\bottomrule
\endfoot
Uses KNIME & 75 & 100.0\% \\
Workflow or nodes described in text & 68 & 90.7\% \\
Workflow screenshots or figures & 48 & 64.0\% \\
Reports KNIME version & 18 & 24.0\% \\
Reports extension/plugin dependencies & 22 & 29.3\% \\
Reports extension installation source & 4 & 5.3\% \\
Reports downloadable workflows & 22 & 29.3\% \\
Articles with successfully downloaded workflows & 18 & 24.0\% \\
\end{longtable}
\endgroup

For workflow retrieval, the authoritative data source is the workflow inventory
derived from the compact audit report and the local workflow directories.
Table~\ref{tab:workflow-retrieval-results} summarizes the recorded
artifact-retrieval outcomes used for manuscript inspection. In this table, the
\emph{Success} column marks whether a KNIME workflow artifact was actually
obtained: \textbf{X} means that a workflow file, archive, or workflow directory
was retrieved, while \textbf{-} means that no KNIME workflow artifact was
obtained or confirmed. Failed or incomplete retrievals are recorded separately
from article-level claims that a workflow was reported or linked.

Table~\ref{tab:workflow-opening-results} then reports the manual opening and
execution checks in the current KNIME environment for the downloaded article
records. Here, the
\emph{Executed} column has a stricter meaning: \textbf{X} means that at least
one workflow for the article opened and executed successfully in the tested
environment, while \textbf{-} means that the workflow opened but did not have a
confirmed successful end-to-end run, or failed during execution. Five article
records had at least one workflow execute successfully in the recorded manual
checks: PAINS, Webinar Pricing Analytics, ImageJ ecosystem integration,
high-content organelle trafficking, and GediNET
\cite{Saubern2011PAINS,VillarroelOrdenes2021KNIMEHub,Dietz2020ImageJKNIME,Pal2018OrganelleKNIME,Qumsiyeh2022GediNET}.
Other opened workflows did not have a confirmed successful end-to-end run in
this manual check. This does not show that those workflows are unreproducible:
they may require additional configuration, dependencies, or repair before a
successful run. The screenshots from the opening checks are archived under
each workflow directory's \texttt{opened/} subdirectory in
\path{data/original/workflows/}.

\begingroup
\scriptsize
\renewcommand{\arraystretch}{1.18}
\begin{longtable}{|%
  >{\centering\arraybackslash}p{0.04\textwidth}|
  >{\raggedright\arraybackslash}p{0.23\textwidth}|
  >{\centering\arraybackslash}p{0.09\textwidth}|
  >{\raggedright\arraybackslash}p{0.54\textwidth}|}
\caption{Workflow retrieval results for article records with attempted or
recorded artifact retrieval outcomes.}
\label{tab:workflow-retrieval-results}\\
\toprule
\# & Article record & Success & Retrieval note \\
\midrule
\endfirsthead
\toprule
\# & Article record & Success & Retrieval note \\
\midrule
\endhead
\bottomrule
\endfoot
1 &
PAINS workflows, DOI \path{10.1002/minf.201100076} &
\textbf{X} &
Automated myExperiment page and download requests returned HTTP 403 HTML, but
we manually obtained both RDKit and Indigo workflow ZIP archives. Both
contain \texttt{workflow.knime} files. \\
2 &
Chemical data curation workflow, DOI \path{10.1186/s13321-018-0315-6} &
\textbf{X} &
GitHub master archive was downloaded after retry and passed ZIP integrity
checking. KNIME workflow files were found in the archive. \\
3 &
ASD genes prediction workflow, DOI \path{10.1371/journal.pone.0208626} &
\textbf{X} &
GitHub archive passed ZIP integrity checking. An inner KNIME workflow archive
was extracted to a workflow directory. \\
4 &
Bitterness/sweetness classifier workflow, DOI \path{10.3389/fchem.2018.00093} &
\textbf{-} &
The original project page redirected to VirtualTaste, and the redirected page
returned HTTP 404. No workflow archive was retrieved. \\
5 &
Medicinal chemistry filters workflow, DOI \path{10.3390/encyclopedia3020035} &
\textbf{X} &
GitLab archive passed ZIP integrity checking and contained KNIME workflow
files. One KNIME Hub URL returned a page, while an older generic Hub URL
returned 404 JSON. \\
6 &
MS-ready structures repository, DOI \path{10.1186/s13321-018-0299-2} &
\textbf{-} &
The GitHub archive passed ZIP integrity checking, but the extracted repository
contained only \texttt{LICENSE} and \texttt{README.md}. No \texttt{.knwf} or
\texttt{workflow.knime} file was found. \\
7 &
KNIME-OMICS reference, DOI \path{10.1016/j.jprot.2016.08.002} &
\textbf{-} &
The referenced GitHub organization page returned HTTP 404. No workflow archive
was retrieved. \\
8 &
QSAR tools workflow, DOI \path{10.1080/07391102.2018.1456975} &
\textbf{-} &
The referenced \path{teqip.jdvu.ac.in/QSAR_Tools/} URL failed DNS resolution.
No workflow archive was retrieved. \\
9 &
Webinar Pricing Analytics workflow, DOI \path{10.1016/j.jbusres.2021.08.036} &
\textbf{X} &
We manually downloaded the KNIME workflow file and added local KNIME
opening screenshots. \\
10 &
Verhaar scheme workflow, DOI \path{10.1016/j.chemosphere.2015.06.009} &
\textbf{-} &
The article reports a supplementary KNIME workflow. Standard Elsevier
supplement URLs resolved as Excel files, but no direct KNIME workflow archive
URL was confirmed. \\
11 &
Spontaneous tail coiling workflow, DOI \path{10.1016/j.ntt.2020.106918} &
\textbf{X} &
We manually downloaded the KNIME workflow file and added local KNIME
opening screenshots. \\
12 &
Caco-2 permeability prediction workflow, DOI \path{10.3390/pharmaceutics14101998} &
\textbf{X} &
We manually downloaded the KNIME workflow file and added local KNIME
opening screenshots. \\
13 &
KniMet workflow, DOI \path{10.1007/s11306-018-1349-5} &
\textbf{X} &
The Zenodo DOI resolved to record 1196407. \texttt{KniMet-master.zip} was
downloaded and contains \texttt{KniMet.knwf}. \\
14 &
NGS sample workflow, DOI \path{10.1093/bioinformatics/btr478} &
\textbf{X} &
We manually downloaded the myExperiment archive after automated requests
were blocked by HTTP 403. The ZIP contains a KNIME workflow directory with
saved data. \\
15 &
ImageJ ecosystem workflows, DOI \path{10.3389/fcomp.2020.00008} &
\textbf{X} &
Six KNIME Hub artifact metadata endpoints returned signed URLs and the
corresponding \texttt{.knwf} files were downloaded. One recorded reference
appears stale or unavailable. \\
16 &
Multi-template matching repository, DOI \path{10.1186/s12859-020-3363-7} &
\textbf{-} &
The GitHub archive was downloaded, but it contained documentation/images only
and no \texttt{.knwf}, \texttt{workflow.knime}, or KNIME workflow archive. A
guessed NodePit URL returned HTTP 404. \\
17 &
GediNET workflows, DOI \path{10.1038/s41598-022-24421-0} &
\textbf{X} &
GitHub master archive was downloaded and contains two \texttt{.knwf} files.
The recorded KNIME Hub short link could not be retrieved. \\
18 &
Nuclear receptor ligand workflow, DOI \path{10.1002/minf.201400188} &
\textbf{-} &
The article reports supplementary workflow material, but Wiley article and
guessed supplementary download endpoints returned Cloudflare challenge pages.
Crossref metadata did not expose a supplementary workflow URL. \\
19 &
High-content organelle trafficking workflow, DOI \path{10.1038/sdata.2018.241} &
\textbf{X} &
The figshare collection DOI resolved through the figshare API. The default
KNIME workflow file was downloaded as \texttt{KNIME workflow.zip}. \\
20 &
Knime4Bio NOTCH2 workflow, DOI \path{10.1093/bioinformatics/btr554} &
\textbf{X} &
The article reports that the workflow was posted on myExperiment. The local
manual artifact is \texttt{The\_NOTCH2\_workflow-v1.zip}, which contains a
KNIME workflow directory. \\
\end{longtable}
\endgroup

\begingroup
\scriptsize
\renewcommand{\arraystretch}{1.18}
\begin{longtable}{|%
  >{\centering\arraybackslash}p{0.04\textwidth}|
  >{\raggedright\arraybackslash}p{0.25\textwidth}|
  >{\centering\arraybackslash}p{0.09\textwidth}|
  >{\raggedright\arraybackslash}p{0.51\textwidth}|}
\caption{Manual opening and execution results for downloaded KNIME workflow
artifacts.}
\label{tab:workflow-opening-results}\\
\toprule
\# & Article record & Executed & Manual KNIME result \\
\midrule
\endfirsthead
\toprule
\# & Article record & Executed & Manual KNIME result \\
\midrule
\endhead
\bottomrule
\endfoot
1 &
PAINS workflows, DOI \path{10.1002/minf.201100076} &
\textbf{X} &
RDKit and Indigo workflow ZIPs were tested. At article-record level, at least
one workflow opened and executed successfully. The separate Indigo workflow
opened but was blocked by a missing Indigo KNIME Integration node. \\
2 &
Chemical data curation workflow, DOI \path{10.1186/s13321-018-0315-6} &
\textbf{-} &
The CompTox rebuilt workflow file opened, but execution failed inside a
metanode during the manual test. \\
3 &
ASD genes prediction workflow, DOI \path{10.1371/journal.pone.0208626} &
\textbf{-} &
The extracted ASD genes prediction workflow opened, but execution failed in an
R node because required R packages and functions were unavailable in the test
environment. \\
4 &
Medicinal chemistry filters workflow, DOI \path{10.3390/encyclopedia3020035} &
\textbf{-} &
The medicinal chemistry filters workflow file opened, but a successful
end-to-end execution was not established in the available notes. \\
5 &
Webinar Pricing Analytics workflow, DOI \path{10.1016/j.jbusres.2021.08.036} &
\textbf{X} &
The Webinar Pricing Analytics workflow file opened and executed successfully
in the current local KNIME environment. \\
6 &
Spontaneous tail coiling workflow, DOI \path{10.1016/j.ntt.2020.106918} &
\textbf{-} &
The spontaneous tail coiling measurement workflow file opened in the current
local KNIME environment, but a successful end-to-end execution was not
established in the available notes. \\
7 &
Caco-2 permeability prediction workflow, DOI \path{10.3390/pharmaceutics14101998} &
\textbf{-} &
The manual Caco-2 prediction \texttt{.knwf} file opened in the current local
KNIME environment, but a successful end-to-end execution was not established
in the available notes. \\
8 &
KniMet workflow, DOI \path{10.1007/s11306-018-1349-5} &
\textbf{-} &
The KniMet workflow from the Zenodo archive opened in the current local KNIME
environment, but successful execution was not established in the available
notes. \\
9 &
NGS sample workflow, DOI \path{10.1093/bioinformatics/btr478} &
\textbf{-} &
The NGS sample workflow archive with saved data opened in the current local
KNIME environment, but successful execution was not established in the
available notes. \\
10 &
ImageJ ecosystem workflows, DOI \path{10.3389/fcomp.2020.00008} &
\textbf{X} &
Six downloaded KNIME Hub \texttt{.knwf} files opened and executed normally in
our test environment after installing additional required KNIME/ImageJ-related
plugins. \\
11 &
GediNET workflows, DOI \path{10.1038/s41598-022-24421-0} &
\textbf{X} &
The GediNET GitHub archive opened in the current local KNIME environment and
the workflow was run successfully in the manual check. \\
12 &
High-content organelle trafficking workflow, DOI \path{10.1038/sdata.2018.241} &
\textbf{X} &
\texttt{KNIME workflow.zip} opened and executed normally in our test
environment after installing additional required plugins. \\
13 &
Knime4Bio NOTCH2 workflow, DOI \path{10.1093/bioinformatics/btr554} &
\textbf{-} &
\texttt{The\_NOTCH2\_workflow-v1.zip} opened in the current local KNIME
environment, but a successful end-to-end execution was not established in the
available notes. \\
\end{longtable}
\endgroup

Taken together, the results show three stages of evidence. First, KNIME is
visible in a broad and sustained scientific literature. Second, reporting
signals are uneven: 68 records describe workflow steps or nodes in text, 49
show workflow screenshots or figures, 22 report downloadable or linked
workflows, and 18 yielded workflow artifacts or directories. Third, workflows
from all 18 records with obtained artifacts opened in the current local KNIME
environment, but only five article records had at least one workflow execute
successfully in the recorded checks. These findings do not show that KNIME
workflows generally fail. They identify a narrower reproducibility risk:
workflow images and textual descriptions do not provide the executable and
environmental evidence needed to test, diagnose, and reproduce the reported
analysis.

\section{Discussion}

This study contributes a semi-automated evidence pipeline and a focused
empirical account of how KNIME workflows are reported and preserved in
scientific publications. Its scope is intentionally narrower than a full
reproducibility benchmark: we do not estimate the failure rate of all published
KNIME workflows. Instead, the pipeline links bibliometric records,
publication-level reporting signals, referenced web resources, retrieved
artifacts, and manual opening and execution outcomes. The resulting contrast is
clear: 68 records contain textual workflow descriptions, 22 report downloadable
or linked workflows, 18 yielded artifacts or directories, and five had a
confirmed successful execution.

The results suggest three practical implications. First, a published workflow
image is not an executable artifact. Several highly cited papers provide
figures or textual descriptions of connected KNIME nodes, but this is not
enough to recover node settings, extension versions, or input paths. Second,
reporting a KNIME version is useful but incomplete.
Workflows may also depend on third-party extension update sites and node
libraries, as shown by the PAINS case where the RDKit workflow could be
updated and executed while the Indigo workflow was blocked by a missing
extension. Third, a reported or linked workflow should not be counted as an
obtained artifact without checking the target resource. The difference between
22 records reporting downloadable or linked workflows and 18 records yielding
workflow artifacts or directories demonstrates the importance of preserving
retrieval evidence separately from publication-level reporting.

The study has several limitations. The OpenAlex query identifies records whose
titles or abstracts match \texttt{KNIME}; it is not a complete census of every
study that used the platform. The top-cited audit is designed to examine
influential records and is not a random sample. Local full text was available
for 79 of the 100 records, and inaccessible records remain in the denominator.
The audit records whether articles report data and code availability but does
not verify every such resource. Finally, the opening and execution checks
reflect the
recorded local environment and available configuration effort. They support
case-level outcomes, not a general estimate of workflow executability.

The next development of this study should move from paper-level evidence to
workflow-level evidence. One direction is to create a larger corpus of public
KNIME workflows from supplements, myExperiment, KNIME Hub, GitHub, and other
repositories, then test whether they open and execute across documented KNIME
versions and extension configurations. A second direction is to preserve
reproducibility environments for selected case studies by recording operating
system, KNIME version, installed
extensions, update sites, input data, import warnings, execution errors, and
output differences.

Tools and services outside the present evidence base could provide a further
layer of usage evidence. \texttt{knime2py} could help identify KNIME workflows
that users try to translate or operationalize outside KNIME
\cite{KaplanKnime2py2026}, while \texttt{k2pweb.org} could provide anonymized,
aggregate evidence about uploaded or executed workflows \cite{K2PWeb2026}.
These sources could help answer a question that the current study leaves open:
which access, dependency, configuration, and execution barriers researchers
encounter when they attempt to reuse KNIME workflows in practice.

\section{Conclusion}

The bibliometric and audit evidence distinguishes the visibility of KNIME in
scientific literature from the preservation of executable KNIME analyses. Of
100 top-cited KNIME-matching records, 75 were classified as KNIME-use cases, 22
reported downloadable or linked workflows, 18 yielded workflow artifacts or
workflow directories, and five had at least one workflow execute successfully
in the recorded checks. These results do not imply that KNIME workflows are
generally unreproducible. They show that long-term assessment depends on
publishing the executable workflow together with version, extension, data, and
execution context, and on recording whether linked artifacts remain retrievable
and runnable.

\bibliographystyle{splncs04}
\bibliography{references}

\end{document}
